\section{Introduction}


Neural networks have become the workhorse non-linear function approximator in a number of machine learning domains, most notably leading to significant advances in reinforcement learning (RL) and supervised learning. 
%
However, it is often overlooked that the success of a particular neural network model depends upon the joint tuning of the model structure, the data presented, and the details of how the model is optimised. Each of these components of a learning framework is controlled by a number of parameters, \emph{hyperparameters}, which influence the learning process and must be properly tuned to fully unlock the network performance.
%
This essential tuning process is computationally expensive, and as neural network systems become more complicated and endowed with even more hyperparameters, the burden of this search process becomes increasingly heavy. 
Furthermore, there is yet another level of complexity in scenarios such as deep reinforcement learning, where the learning problem itself can be highly non-stationary (\eg~dependent on which parts of an environment an agent is currently able to explore). As a consequence, it might be the case that the ideal hyperparameters for such learning problems are themselves highly non-stationary, and should vary in a way that precludes setting their schedule in advance.

Two common tracks for the tuning of hyperparameters exist: \emph{parallel search} and \emph{sequential optimisation}, which trade-off concurrently used computational resources with the time required to achieve optimal results. 
%
Parallel search performs many parallel optimisation processes (by \emph{optimisation process} we refer to neural network training runs), each with different hyperparameters, with a view to finding a single best output from one of the optimisation processes -- examples of this are grid search and random search. 
%
Sequential optimisation performs few optimisation processes in parallel, but does so many times sequentially, to gradually perform hyperparameter optimisation using information obtained from earlier training runs to inform later ones -- examples of this are hand tuning and Bayesian optimisation. Sequential optimisation will in general provide the best solutions, but requires multiple sequential training runs, which is often unfeasible for lengthy optimisation processes.


In this work, we present a simple method, Population Based Training (PBT) which bridges and extends parallel search methods and sequential optimisation methods.
%
Advantageously, our proposal has a wall-clock run time that is no greater than that of a single optimisation process, does not require sequential runs, and is also able to use fewer computational resources than naive search methods such as random or grid search.
%
Our approach leverages information sharing across a population of concurrently running optimisation processes, and allows for online propagation/transfer of parameters and hyperparameters between members of the population based on their performance. Furthermore, unlike most other adaptation schemes, our method is capable of performing online adaptation of hyperparameters -- which can be particularly important in problems with highly non-stationary learning dynamics, such as reinforcement learning settings.
%
PBT is decentralised and asynchronous, requiring minimal overhead and infrastructure. While inherently greedy, we show that this meta-optimisation process results in effective and automatic tuning of hyperparameters, allowing them to be adaptive throughout training. In addition, the model selection and propagation process ensures that intermediate good models are given more computational resources, and are used as a basis of further optimisation and hyperparameter search. 

We apply PBT to a diverse set of problems and domains to demonstrate its effectiveness and wide applicability. Firstly, we look at the problem of deep reinforcement learning, showing how PBT can be used to optimise UNREAL~\citep{jaderberg2016reinforcement} on DeepMind Lab levels~\citep{beattie2016deepmind}, Feudal Networks~\citep{vezhnevets2017feudal} on Atari games~\citep{bellemare2013arcade}, and simple A3C agents for StarCraft II~\citep{vinyals2017starcraft}. In all three cases we show faster learning and higher performance across a suite of tasks, with PBT allowing discovery of new state-of-the-art performance and behaviour, as well as the potential to reduce the computational resources for training. Secondly, we look at its use in supervised learning for machine translation on WMT 2014 English-to-German with Transformer networks~\citep{vaswani2017attention}, showing that PBT can match and even outperform heavily tuned hyperparameter schedules by optimising for BLEU score directly. Finally, we apply PBT to the training of Generative Adversarial Networks (GANs)~\citep{goodfellowgan} by optimising for the Inception score~\citep{openaigan} -- the result is stable GAN training with
large performance improvements over a strong, well-tuned baseline with the same architecture~\citep{dcgan}.

The improvements we show empirically are the result of (a) automatic selection of hyperparameters during training,  (b) online model selection to maximise the use of computation spent on promising models, and (c) the ability for online adaptation of hyperparameters to enable non-stationary training regimes and the discovery of complex hyperparameter schedules. 

In \sref{sec:related} we review related work on parallel and sequential hyperparameter optimisation techniques, as well as evolutionary optimisation methods which bear a resemblance to PBT. We introduce PBT in \sref{sec:pbt}, outlining the algorithm in a general manner which is used as a basis for experiments. The specific incarnations of PBT and the results of experiments in different domains are given in \sref{sec:exp}, and finally we conclude in \sref{sec:conclusion}.

\begin{figure}[t]
\centering
\includegraphics[width=1.0\textwidth]{figures/overview.pdf}
\caption{\small Common paradigms of tuning hyperparameters: \emph{sequential optimisation} and \emph{parallel search}, as compared to the method of \emph{population based training} introduced in this work. (a) Sequential optimisation requires multiple training runs to be completed (potentially with early stopping), after which new hyperparameters are selected and the model is retrained from scratch with the new hyperparameters. This is an inherently sequential process and leads to long hyperparameter optimisation times, though uses minimal computational resources. (b) Parallel random/grid search of hyperparameters trains multiple models in parallel with different weight initialisations and hyperparameters, with the view that one of the models will be optimised the best. This only requires the time for one training run, but requires the use of more computational resources to train many models in parallel. (c) Population based training starts like parallel search, randomly sampling hyperparameters and weight initialisations. However, each training run asynchronously evaluates its performance periodically. If a model in the population is under-performing, it will \emph{exploit} the rest of the population by replacing itself with a better performing model, and it will \emph{explore} new hyperparameters by modifying the better model's hyperparameters, before training is continued. This process allows hyperparameters to be optimised online, and the computational resources to be focused on the hyperparameter and weight space that has most chance of producing good results. The result is a hyperparameter tuning method that while very simple, results in faster learning, lower computational resources, and often better solutions.}
\label{fig:overview}
\end{figure}
